% Created 2026-07-15 mer 14:00
% Intended LaTeX compiler: lualatex
\documentclass[12pt,letterpaper]{article}
\usepackage{geometry}
\geometry{margin=3cm}
\usepackage{amssymb}
\setcounter{secnumdepth}{0}
\usepackage[french]{babel}
\usepackage{amsmath}
\usepackage{fontspec}
\usepackage{graphicx}
\usepackage{longtable}
\usepackage{wrapfig}
\usepackage{rotating}
\usepackage[normalem]{ulem}
\usepackage{capt-of}
\usepackage{hyperref}
\date{}
\title{Raccourcis de DWL}
\hypersetup{
 pdfauthor={},
 pdftitle={Raccourcis de DWL},
 pdfkeywords={},
 pdfsubject={},
 pdfcreator={},
 pdflang={French}}
\begin{document}

\maketitle
\begin{center}
\begin{tabular}{ll}
Raccourcis & action\\
\hline
Super + p & password menu\\
Super + a & application menu\\
Super + s & system menu\\
Super + z & zwift\\
Super + Shift + Return & terminal\\
Super + o & vimwiki\\
Super + n & firefox\\
Super + r & focus stack down\\
Super + t & focus stack up\\
Super + g & incnmaster +1 augmente le nb de fenêtre dans le maître\\
Super + , & incnmaster -1 diminue le nb de fenêtre dans le maître\\
Super + l & déplace la séparation vers la gauche\\
Super + i & déplace la séparation vers la droite\\
Super + Return & zoom : place la fenêtre active première dans le master\\
Super + Tab & view : revient au Tag précédent\\
Super + Shift + c & fermer la fenêtre\\
Super + Shift + q & fermer la session\\
Super + Shift + w & tile mode\\
Super + Shift + f & float mode\\
Super + Shift + b & monocle mode\\
Super + space & placer la fenêtre active en floating\\
Super + Shift + f & bascule floating\\
Super + e & bascule en plein écran\\
Super + 0 & montre toutes les fenêtres dans le même écran\\
Super + @ & épingle la fenêtre active sur tous les écrans\\
Super + m & focus écran de gauche\\
Super + d & focus écran de droite\\
Super + Shift + m & déplace sur l'écran de gauche\\
Super + Shift + d & déplace sur l'écran de droite\\
 & \\
\end{tabular}
\end{center}
\end{document}
